\documentclass[1pt]{article}

\usepackage[breakable]{tcolorbox}
\usepackage{parskip} % Stop auto-indenting (to mimic markdown behaviour)


% Basic figure setup, for now with no caption control since it's done
% automatically by Pandoc (which extracts ![](path) syntax from Markdown).
\usepackage{graphicx}
% Keep aspect ratio if custom image width or height is specified
\setkeys{Gin}{keepaspectratio}
% Maintain compatibility with old templates. Remove in nbconvert 6.0
\let\Oldincludegraphics\includegraphics
% Ensure that by default, figures have no caption (until we provide a
% proper Figure object with a Caption API and a way to capture that
% in the conversion process - todo).
\usepackage{caption}
\DeclareCaptionFormat{nocaption}{}
\captionsetup{format=nocaption,aboveskip=0pt,belowskip=0pt}

\usepackage{float}
\floatplacement{figure}{H} % forces figures to be placed at the correct location
\usepackage{xcolor} % Allow colors to be defined
\usepackage{enumerate} % Needed for markdown enumerations to work
\usepackage{geometry} % Used to adjust the document margins
\usepackage{amsmath} % Equations
\usepackage{amssymb} % Equations
\usepackage{textcomp} % defines textquotesingle
% Hack from http://tex.stackexchange.com/a/47451/13684:
\AtBeginDocument{%
    \def\PYZsq{\textquotesingle}% Upright quotes in Pygmentized code
}
\usepackage{upquote} % Upright quotes for verbatim code
\usepackage{eurosym} % defines \euro

\usepackage{iftex}
\ifPDFTeX
\usepackage[T1]{fontenc}
\IfFileExists{alphabeta.sty}{
    \usepackage{alphabeta}
}{
    \usepackage[mathletters]{ucs}
    \usepackage[utf8x]{inputenc}
}
\else
\usepackage{fontspec}
\usepackage{unicode-math}
\setmainfont{JetBrains Mono}
\fi

\usepackage{fancyhdr}

\pagestyle{fancy}
\fancyhf{}

\fancyfoot[R]{%
    \parbox[b]{5cm}{%
        \normalfont\normalsize
        \raggedright
        Name: Swapnaraj Mohanty\\
        SIC No: 24BCSH93\\
        Lab Roll No: 11\\
        Date: 17$^{\mathrm{th}}$ August 2026
    }%
}
\renewcommand{\headrulewidth}{0pt}
\renewcommand{\footrulewidth}{0pt}

\usepackage{fancyvrb} % verbatim replacement that allows latex
\usepackage{grffile} % extends the file name processing of package graphics
% to support a larger range
\makeatletter % fix for old versions of grffile with XeLaTeX
\@ifpackagelater{grffile}{2019/11/01}
{
% Do nothing on new versions
}
{
    \def\Gread@@xetex#1{%
    \IfFileExists{"\Gin@base".bb}%
    {\Gread@eps{\Gin@base.bb}}%
    {\Gread@@xetex@aux#1}%
}
}
\makeatother
\usepackage[Export]{adjustbox} % Used to constrain images to a maximum size
\adjustboxset{max size={0.9\linewidth}{0.9\paperheight}}

% The hyperref package gives us a pdf with properly built
% internal navigation ('pdf bookmarks' for the table of contents,
% internal cross-reference links, web links for URLs, etc.)
\usepackage{hyperref}
% The default LaTeX title has an obnoxious amount of whitespace. By default,
% titling removes some of it. It also provides customization options.
\usepackage{titling}
\usepackage{longtable} % longtable support required by pandoc >1.10
\usepackage{booktabs}  % table support for pandoc > 1.12.2
\usepackage{array}     % table support for pandoc >= 2.11.3
\usepackage{calc}      % table minipage width calculation for pandoc >= 2.11.1
\usepackage[inline]{enumitem} % IRkernel/repr support (it uses the enumerate* environment)
\usepackage[normalem]{ulem} % ulem is needed to support strikethroughs (\sout)
% normalem makes italics be italics, not underlines
\usepackage{soul}      % strikethrough (\st) support for pandoc >= 3.0.0
\usepackage{mathrsfs}


% Colors for the hyperref package
\definecolor{urlcolor}{rgb}{0,.145,.698}
\definecolor{linkcolor}{rgb}{.71,0.21,0.01}
\definecolor{citecolor}{rgb}{.12,.54,.11}

% ANSI colors
\definecolor{ansi-black}{HTML}{3E424D}
\definecolor{ansi-black-intense}{HTML}{282C36}
\definecolor{ansi-red}{HTML}{E75C58}
\definecolor{ansi-red-intense}{HTML}{B22B31}
\definecolor{ansi-green}{HTML}{00A250}
\definecolor{ansi-green-intense}{HTML}{007427}
\definecolor{ansi-yellow}{HTML}{DDB62B}
\definecolor{ansi-yellow-intense}{HTML}{B27D12}
\definecolor{ansi-blue}{HTML}{208FFB}
\definecolor{ansi-blue-intense}{HTML}{0065CA}
\definecolor{ansi-magenta}{HTML}{D160C4}
\definecolor{ansi-magenta-intense}{HTML}{A03196}
\definecolor{ansi-cyan}{HTML}{60C6C8}
\definecolor{ansi-cyan-intense}{HTML}{258F8F}
\definecolor{ansi-white}{HTML}{C5C1B4}
\definecolor{ansi-white-intense}{HTML}{A1A6B2}
\definecolor{ansi-default-inverse-fg}{HTML}{FFFFFF}
\definecolor{ansi-default-inverse-bg}{HTML}{000000}

% common color for the border for error outputs.
\definecolor{outerrorbackground}{HTML}{FFDFDF}

% commands and environments needed by pandoc snippets
% extracted from the output of `pandoc -s`
\providecommand{\tightlist}{%
    \setlength{\itemsep}{0pt}\setlength{\parskip}{0pt}}
\DefineVerbatimEnvironment{Highlighting}{Verbatim}{commandchars=\\\{\}}
% Add ',fontsize=\small' for more characters per line
\newenvironment{Shaded}{}{}
\newcommand{\KeywordTok}[1]{\textcolor[rgb]{0.00,0.44,0.13}{\textbf{{#1}}}}
\newcommand{\DataTypeTok}[1]{\textcolor[rgb]{0.56,0.13,0.00}{{#1}}}
\newcommand{\DecValTok}[1]{\textcolor[rgb]{0.25,0.63,0.44}{{#1}}}
\newcommand{\BaseNTok}[1]{\textcolor[rgb]{0.25,0.63,0.44}{{#1}}}
\newcommand{\FloatTok}[1]{\textcolor[rgb]{0.25,0.63,0.44}{{#1}}}
\newcommand{\CharTok}[1]{\textcolor[rgb]{0.25,0.44,0.63}{{#1}}}
\newcommand{\StringTok}[1]{\textcolor[rgb]{0.25,0.44,0.63}{{#1}}}
\newcommand{\CommentTok}[1]{\textcolor[rgb]{0.38,0.63,0.69}{\textit{{#1}}}}
\newcommand{\OtherTok}[1]{\textcolor[rgb]{0.00,0.44,0.13}{{#1}}}
\newcommand{\AlertTok}[1]{\textcolor[rgb]{1.00,0.00,0.00}{\textbf{{#1}}}}
\newcommand{\FunctionTok}[1]{\textcolor[rgb]{0.02,0.16,0.49}{{#1}}}
\newcommand{\RegionMarkerTok}[1]{{#1}}
\newcommand{\ErrorTok}[1]{\textcolor[rgb]{1.00,0.00,0.00}{\textbf{{#1}}}}
\newcommand{\NormalTok}[1]{{#1}}

% Additional commands for more recent versions of Pandoc
\newcommand{\ConstantTok}[1]{\textcolor[rgb]{0.53,0.00,0.00}{{#1}}}
\newcommand{\SpecialCharTok}[1]{\textcolor[rgb]{0.25,0.44,0.63}{{#1}}}
\newcommand{\VerbatimStringTok}[1]{\textcolor[rgb]{0.25,0.44,0.63}{{#1}}}
\newcommand{\SpecialStringTok}[1]{\textcolor[rgb]{0.73,0.40,0.53}{{#1}}}
\newcommand{\ImportTok}[1]{{#1}}
\newcommand{\DocumentationTok}[1]{\textcolor[rgb]{0.73,0.13,0.13}{\textit{{#1}}}}
\newcommand{\AnnotationTok}[1]{\textcolor[rgb]{0.38,0.63,0.69}{\textbf{\textit{{#1}}}}}
\newcommand{\CommentVarTok}[1]{\textcolor[rgb]{0.38,0.63,0.69}{\textbf{\textit{{#1}}}}}
\newcommand{\VariableTok}[1]{\textcolor[rgb]{0.10,0.09,0.49}{{#1}}}
\newcommand{\ControlFlowTok}[1]{\textcolor[rgb]{0.00,0.44,0.13}{\textbf{{#1}}}}
\newcommand{\OperatorTok}[1]{\textcolor[rgb]{0.40,0.40,0.40}{{#1}}}
\newcommand{\BuiltInTok}[1]{{#1}}
\newcommand{\ExtensionTok}[1]{{#1}}
\newcommand{\PreprocessorTok}[1]{\textcolor[rgb]{0.74,0.48,0.00}{{#1}}}
\newcommand{\AttributeTok}[1]{\textcolor[rgb]{0.49,0.56,0.16}{{#1}}}
\newcommand{\InformationTok}[1]{\textcolor[rgb]{0.38,0.63,0.69}{\textbf{\textit{{#1}}}}}
\newcommand{\WarningTok}[1]{\textcolor[rgb]{0.38,0.63,0.69}{\textbf{\textit{{#1}}}}}
\makeatletter
\newsavebox\pandoc@box
\newcommand*\pandocbounded[1]{%
    \sbox\pandoc@box{#1}%
    % scaling factors for width and height
    \Gscale@div\@tempa\textheight{\dimexpr\ht\pandoc@box+\dp\pandoc@box\relax}%
    \Gscale@div\@tempb\linewidth{\wd\pandoc@box}%
% select the smaller of both
    \ifdim
        \@tempb\p@<\@tempa\p@
        \let\@tempa\@tempb
    \fi
% scaling accordingly (\@tempa < 1)
    \ifdim
        \@tempa\p@<\p@
        \scalebox{\@tempa}{\usebox\pandoc@box}%
    % scaling not needed, use as it is
    \else
        \usebox{\pandoc@box}%
    \fi
}
\makeatother

% Define a nice break command that doesn't care if a line doesn't already
% exist.
\def\br{\hspace*{\fill} \\* }
% Math Jax compatibility definitions
\def\gt{>}
\def\lt{<}
\let\Oldtex\TeX
\let\Oldlatex\LaTeX
\renewcommand{\TeX}{\textrm{\Oldtex}}
\renewcommand{\LaTeX}{\textrm{\Oldlatex}}
% Document parameters
% Document title
\title{day2}


% Pygments definitions
\makeatletter
\def\PY@reset{\let\PY@it=\relax \let\PY@bf=\relax%
\let\PY@ul=\relax \let\PY@tc=\relax%
\let\PY@bc=\relax \let\PY@ff=\relax}
\def\PY@tok#1{\csname PY@tok@#1\endcsname}
\def\PY@toks#1+{\ifx
                    \relax#1\empty\else%
                    \PY@tok{#1}\expandafter\PY@toks
\fi}
\def\PY@do#1{\PY@bc{\PY@tc{\PY@ul{%
    \PY@it{\PY@bf{\PY@ff{#1}}}}}}}
\def\PY#1#2{\PY@reset\PY@toks#1+\relax+\PY@do{#2}}

\@namedef{PY@tok@w}{\def\PY@tc##1{\textcolor[rgb]{0.73,0.73,0.73}{##1}}}
\@namedef{PY@tok@c}{\let\PY@it=\textit\def\PY@tc##1{\textcolor[rgb]{0.24,0.48,0.48}{##1}}}
\@namedef{PY@tok@cp}{\def\PY@tc##1{\textcolor[rgb]{0.61,0.40,0.00}{##1}}}
\@namedef{PY@tok@k}{\let\PY@bf=\textbf\def\PY@tc##1{\textcolor[rgb]{0.00,0.50,0.00}{##1}}}
\@namedef{PY@tok@kp}{\def\PY@tc##1{\textcolor[rgb]{0.00,0.50,0.00}{##1}}}
\@namedef{PY@tok@kt}{\def\PY@tc##1{\textcolor[rgb]{0.69,0.00,0.25}{##1}}}
\@namedef{PY@tok@o}{\def\PY@tc##1{\textcolor[rgb]{0.40,0.40,0.40}{##1}}}
\@namedef{PY@tok@ow}{\let\PY@bf=\textbf\def\PY@tc##1{\textcolor[rgb]{0.67,0.13,1.00}{##1}}}
\@namedef{PY@tok@nb}{\def\PY@tc##1{\textcolor[rgb]{0.00,0.50,0.00}{##1}}}
\@namedef{PY@tok@nf}{\def\PY@tc##1{\textcolor[rgb]{0.00,0.00,1.00}{##1}}}
\@namedef{PY@tok@nc}{\let\PY@bf=\textbf\def\PY@tc##1{\textcolor[rgb]{0.00,0.00,1.00}{##1}}}
\@namedef{PY@tok@nn}{\let\PY@bf=\textbf\def\PY@tc##1{\textcolor[rgb]{0.00,0.00,1.00}{##1}}}
\@namedef{PY@tok@ne}{\let\PY@bf=\textbf\def\PY@tc##1{\textcolor[rgb]{0.80,0.25,0.22}{##1}}}
\@namedef{PY@tok@nv}{\def\PY@tc##1{\textcolor[rgb]{0.10,0.09,0.49}{##1}}}
\@namedef{PY@tok@no}{\def\PY@tc##1{\textcolor[rgb]{0.53,0.00,0.00}{##1}}}
\@namedef{PY@tok@nl}{\def\PY@tc##1{\textcolor[rgb]{0.46,0.46,0.00}{##1}}}
\@namedef{PY@tok@ni}{\let\PY@bf=\textbf\def\PY@tc##1{\textcolor[rgb]{0.44,0.44,0.44}{##1}}}
\@namedef{PY@tok@na}{\def\PY@tc##1{\textcolor[rgb]{0.41,0.47,0.13}{##1}}}
\@namedef{PY@tok@nt}{\let\PY@bf=\textbf\def\PY@tc##1{\textcolor[rgb]{0.00,0.50,0.00}{##1}}}
\@namedef{PY@tok@nd}{\def\PY@tc##1{\textcolor[rgb]{0.67,0.13,1.00}{##1}}}
\@namedef{PY@tok@s}{\def\PY@tc##1{\textcolor[rgb]{0.73,0.13,0.13}{##1}}}
\@namedef{PY@tok@sd}{\let\PY@it=\textit\def\PY@tc##1{\textcolor[rgb]{0.73,0.13,0.13}{##1}}}
\@namedef{PY@tok@si}{\let\PY@bf=\textbf\def\PY@tc##1{\textcolor[rgb]{0.64,0.35,0.47}{##1}}}
\@namedef{PY@tok@se}{\let\PY@bf=\textbf\def\PY@tc##1{\textcolor[rgb]{0.67,0.36,0.12}{##1}}}
\@namedef{PY@tok@sr}{\def\PY@tc##1{\textcolor[rgb]{0.64,0.35,0.47}{##1}}}
\@namedef{PY@tok@ss}{\def\PY@tc##1{\textcolor[rgb]{0.10,0.09,0.49}{##1}}}
\@namedef{PY@tok@sx}{\def\PY@tc##1{\textcolor[rgb]{0.00,0.50,0.00}{##1}}}
\@namedef{PY@tok@m}{\def\PY@tc##1{\textcolor[rgb]{0.40,0.40,0.40}{##1}}}
\@namedef{PY@tok@gh}{\let\PY@bf=\textbf\def\PY@tc##1{\textcolor[rgb]{0.00,0.00,0.50}{##1}}}
\@namedef{PY@tok@gu}{\let\PY@bf=\textbf\def\PY@tc##1{\textcolor[rgb]{0.50,0.00,0.50}{##1}}}
\@namedef{PY@tok@gd}{\def\PY@tc##1{\textcolor[rgb]{0.63,0.00,0.00}{##1}}}
\@namedef{PY@tok@gi}{\def\PY@tc##1{\textcolor[rgb]{0.00,0.52,0.00}{##1}}}
\@namedef{PY@tok@gr}{\def\PY@tc##1{\textcolor[rgb]{0.89,0.00,0.00}{##1}}}
\@namedef{PY@tok@ge}{\let\PY@it=\textit}
\@namedef{PY@tok@gs}{\let\PY@bf=\textbf}
\@namedef{PY@tok@ges}{\let\PY@bf=\textbf\let\PY@it=\textit}
\@namedef{PY@tok@gp}{\let\PY@bf=\textbf\def\PY@tc##1{\textcolor[rgb]{0.00,0.00,0.50}{##1}}}
\@namedef{PY@tok@go}{\def\PY@tc##1{\textcolor[rgb]{0.44,0.44,0.44}{##1}}}
\@namedef{PY@tok@gt}{\def\PY@tc##1{\textcolor[rgb]{0.00,0.27,0.87}{##1}}}
\@namedef{PY@tok@err}{\def\PY@bc##1{{\setlength{\fboxsep}{\string -\fboxrule}\fcolorbox[rgb]{1.00,0.00,0.00}{1,1,1}{\strut ##1}}}}
\@namedef{PY@tok@kc}{\let\PY@bf=\textbf\def\PY@tc##1{\textcolor[rgb]{0.00,0.50,0.00}{##1}}}
\@namedef{PY@tok@kd}{\let\PY@bf=\textbf\def\PY@tc##1{\textcolor[rgb]{0.00,0.50,0.00}{##1}}}
\@namedef{PY@tok@kn}{\let\PY@bf=\textbf\def\PY@tc##1{\textcolor[rgb]{0.00,0.50,0.00}{##1}}}
\@namedef{PY@tok@kr}{\let\PY@bf=\textbf\def\PY@tc##1{\textcolor[rgb]{0.00,0.50,0.00}{##1}}}
\@namedef{PY@tok@bp}{\def\PY@tc##1{\textcolor[rgb]{0.00,0.50,0.00}{##1}}}
\@namedef{PY@tok@fm}{\def\PY@tc##1{\textcolor[rgb]{0.00,0.00,1.00}{##1}}}
\@namedef{PY@tok@vc}{\def\PY@tc##1{\textcolor[rgb]{0.10,0.09,0.49}{##1}}}
\@namedef{PY@tok@vg}{\def\PY@tc##1{\textcolor[rgb]{0.10,0.09,0.49}{##1}}}
\@namedef{PY@tok@vi}{\def\PY@tc##1{\textcolor[rgb]{0.10,0.09,0.49}{##1}}}
\@namedef{PY@tok@vm}{\def\PY@tc##1{\textcolor[rgb]{0.10,0.09,0.49}{##1}}}
\@namedef{PY@tok@sa}{\def\PY@tc##1{\textcolor[rgb]{0.73,0.13,0.13}{##1}}}
\@namedef{PY@tok@sb}{\def\PY@tc##1{\textcolor[rgb]{0.73,0.13,0.13}{##1}}}
\@namedef{PY@tok@sc}{\def\PY@tc##1{\textcolor[rgb]{0.73,0.13,0.13}{##1}}}
\@namedef{PY@tok@dl}{\def\PY@tc##1{\textcolor[rgb]{0.73,0.13,0.13}{##1}}}
\@namedef{PY@tok@s2}{\def\PY@tc##1{\textcolor[rgb]{0.73,0.13,0.13}{##1}}}
\@namedef{PY@tok@sh}{\def\PY@tc##1{\textcolor[rgb]{0.73,0.13,0.13}{##1}}}
\@namedef{PY@tok@s1}{\def\PY@tc##1{\textcolor[rgb]{0.73,0.13,0.13}{##1}}}
\@namedef{PY@tok@mb}{\def\PY@tc##1{\textcolor[rgb]{0.40,0.40,0.40}{##1}}}
\@namedef{PY@tok@mf}{\def\PY@tc##1{\textcolor[rgb]{0.40,0.40,0.40}{##1}}}
\@namedef{PY@tok@mh}{\def\PY@tc##1{\textcolor[rgb]{0.40,0.40,0.40}{##1}}}
\@namedef{PY@tok@mi}{\def\PY@tc##1{\textcolor[rgb]{0.40,0.40,0.40}{##1}}}
\@namedef{PY@tok@il}{\def\PY@tc##1{\textcolor[rgb]{0.40,0.40,0.40}{##1}}}
\@namedef{PY@tok@mo}{\def\PY@tc##1{\textcolor[rgb]{0.40,0.40,0.40}{##1}}}
\@namedef{PY@tok@ch}{\let\PY@it=\textit\def\PY@tc##1{\textcolor[rgb]{0.24,0.48,0.48}{##1}}}
\@namedef{PY@tok@cm}{\let\PY@it=\textit\def\PY@tc##1{\textcolor[rgb]{0.24,0.48,0.48}{##1}}}
\@namedef{PY@tok@cpf}{\let\PY@it=\textit\def\PY@tc##1{\textcolor[rgb]{0.24,0.48,0.48}{##1}}}
\@namedef{PY@tok@c1}{\let\PY@it=\textit\def\PY@tc##1{\textcolor[rgb]{0.24,0.48,0.48}{##1}}}
\@namedef{PY@tok@cs}{\let\PY@it=\textit\def\PY@tc##1{\textcolor[rgb]{0.24,0.48,0.48}{##1}}}

\def\PYZbs{\char`\\}
\def\PYZus{\char`\_}
\def\PYZob{\char`\{}
\def\PYZcb{\char`\}}
\def\PYZca{\char`\^}
\def\PYZam{\char`\&}
\def\PYZlt{\char`\<}
\def\PYZgt{\char`\>}
\def\PYZsh{\char`\#}
\def\PYZpc{\char`\%}
\def\PYZdl{\char`\$}
\def\PYZhy{\char`\-}
\def\PYZsq{\char`\'}
\def\PYZdq{\char`\"}
\def\PYZti{\char`\~}
% for compatibility with earlier versions
\def\PYZat{@}
\def\PYZlb{[}
\def\PYZrb{]}
\makeatother


% For linebreaks inside Verbatim environment from package fancyvrb.
\makeatletter
\newbox\Wrappedcontinuationbox
\newbox\Wrappedvisiblespacebox
\newcommand*\Wrappedvisiblespace {\textcolor{red}{\textvisiblespace}}
\newcommand*\Wrappedcontinuationsymbol {\textcolor{red}{\llap{\tiny$\m@th\hookrightarrow$}}}
\newcommand*\Wrappedcontinuationindent {3ex }
\newcommand*\Wrappedafterbreak {\kern\Wrappedcontinuationindent\copy\Wrappedcontinuationbox}
% Take advantage of the already applied Pygments mark-up to insert
% potential linebreaks for TeX processing.
%        {, <, #, %, $, ' and ": go to next line.
%        _, }, ^, &, >, - and ~: stay at end of broken line.
% Use of \textquotesingle for straight quote.
\newcommand*\Wrappedbreaksatspecials {%
    \def\PYGZus{\discretionary{\char`\_}{\Wrappedafterbreak}{\char`\_}}%
    \def\PYGZob{\discretionary{}{\Wrappedafterbreak\char`\{}{\char`\{}}%
    \def\PYGZcb{\discretionary{\char`\}}{\Wrappedafterbreak}{\char`\}}}%
    \def\PYGZca{\discretionary{\char`\^}{\Wrappedafterbreak}{\char`\^}}%
    \def\PYGZam{\discretionary{\char`\&}{\Wrappedafterbreak}{\char`\&}}%
    \def\PYGZlt{\discretionary{}{\Wrappedafterbreak\char`\<}{\char`\<}}%
    \def\PYGZgt{\discretionary{\char`\>}{\Wrappedafterbreak}{\char`\>}}%
    \def\PYGZsh{\discretionary{}{\Wrappedafterbreak\char`\#}{\char`\#}}%
    \def\PYGZpc{\discretionary{}{\Wrappedafterbreak\char`\%}{\char`\%}}%
    \def\PYGZdl{\discretionary{}{\Wrappedafterbreak\char`\$}{\char`\$}}%
    \def\PYGZhy{\discretionary{\char`\-}{\Wrappedafterbreak}{\char`\-}}%
    \def\PYGZsq{\discretionary{}{\Wrappedafterbreak\textquotesingle}{\textquotesingle}}%
    \def\PYGZdq{\discretionary{}{\Wrappedafterbreak\char`\"}{\char`\"}}%
    \def\PYGZti{\discretionary{\char`\~}{\Wrappedafterbreak}{\char`\~}}%
}
% Some characters . , ; ? ! / are not pygmentized.
% This macro makes them "active" and they will insert potential linebreaks
\newcommand*\Wrappedbreaksatpunct {%
    \lccode`\~`\.\lowercase{\def~}{\discretionary{\hbox{\char`\.}}{\Wrappedafterbreak}{\hbox{\char`\.}}}%
    \lccode`\~`\,\lowercase{\def~}{\discretionary{\hbox{\char`\,}}{\Wrappedafterbreak}{\hbox{\char`\,}}}%
    \lccode`\~`\;\lowercase{\def~}{\discretionary{\hbox{\char`\;}}{\Wrappedafterbreak}{\hbox{\char`\;}}}%
    \lccode`\~`\:\lowercase{\def~}{\discretionary{\hbox{\char`\:}}{\Wrappedafterbreak}{\hbox{\char`\:}}}%
    \lccode`\~`\?\lowercase{\def~}{\discretionary{\hbox{\char`\?}}{\Wrappedafterbreak}{\hbox{\char`\?}}}%
    \lccode`\~`\!\lowercase{\def~}{\discretionary{\hbox{\char`\!}}{\Wrappedafterbreak}{\hbox{\char`\!}}}%
    \lccode`\~`\/\lowercase{\def~}{\discretionary{\hbox{\char`\/}}{\Wrappedafterbreak}{\hbox{\char`\/}}}%
    \catcode`\.\active
    \catcode`\,\active
    \catcode`\;\active
    \catcode`\:\active
    \catcode`\?\active
    \catcode`\!\active
    \catcode`\/\active
    \lccode`\~`\~
}
\makeatother

\let\OriginalVerbatim=\Verbatim
\makeatletter
\renewcommand{\Verbatim}[1][1]{%
%\parskip\z@skip
    \sbox\Wrappedcontinuationbox {\Wrappedcontinuationsymbol}%
    \sbox\Wrappedvisiblespacebox {\FV@SetupFont\Wrappedvisiblespace}%
    \def\FancyVerbFormatLine ##1{\hsize\linewidth
    \vtop{\raggedright\hyphenpenalty\z@\exhyphenpenalty\z@
    \doublehyphendemerits\z@\finalhyphendemerits\z@
    \strut ##1\strut}%
    }%
% If the linebreak is at a space, the latter will be displayed as visible
% space at end of first line, and a continuation symbol starts next line.
% Stretch/shrink are however usually zero for typewriter font.
    \def\FV@Space {%
        \nobreak\hskip\z@ plus\fontdimen3\font minus\fontdimen4\font
        \discretionary{\copy\Wrappedvisiblespacebox}{\Wrappedafterbreak}
        {\kern\fontdimen2\font}%
    }%

    % Allow breaks at special characters using \PYG... macros.
    \Wrappedbreaksatspecials
    % Breaks at punctuation characters . , ; ? ! and / need catcode=\active
    \OriginalVerbatim[#1,codes*=\Wrappedbreaksatpunct]%
}
\makeatother

% Exact colors from NB
\definecolor{incolor}{HTML}{303F9F}
\definecolor{outcolor}{HTML}{D84315}
\definecolor{cellborder}{HTML}{CFCFCF}
\definecolor{cellbackground}{HTML}{F7F7F7}

% prompt
\makeatletter
\newcommand{\boxspacing}{\kern\kvtcb@left@rule\kern\kvtcb@boxsep}
\makeatother
\newcommand{\prompt}[4]{
        {\ttfamily\llap{{\color{#2}[#3]:\hspace{3pt}#4}}\vspace{-\baselineskip}}
}


% Prevent overflowing lines due to hard-to-break entities
\sloppy
% Setup hyperref package
\hypersetup{
    breaklinks=true,  % so long urls are correctly broken across lines
    colorlinks=true,
    urlcolor=urlcolor,
    linkcolor=linkcolor,
    citecolor=citecolor,
}
% Slightly bigger margins than the latex defaults

\geometry{verbose,tmargin=1in,bmargin=1in,lmargin=1in,rmargin=1in}



\begin{document}
    \hypertarget{data-pre-processing}{%
        \section{Data Pre-Processing}\label{data-pre-processing}}

    \begin{tcolorbox}[breakable, size=fbox, boxrule=1pt, pad at break*=1mm,colback=cellbackground, colframe=cellborder]
        \prompt{In}{incolor}{1}{\boxspacing}
        \begin{Verbatim}[commandchars=\\\{\}]
\PY{k+kn}{import}\PY{+w}{ }\PY{n+nn}{numpy}\PY{+w}{ }\PY{k}{as}\PY{+w}{ }\PY{n+nn}{np}
\PY{k+kn}{import}\PY{+w}{ }\PY{n+nn}{matplotlib}\PY{n+nn}{.}\PY{n+nn}{pyplot}\PY{+w}{ }\PY{k}{as}\PY{+w}{ }\PY{n+nn}{plt}
\PY{k+kn}{import}\PY{+w}{ }\PY{n+nn}{pandas}\PY{+w}{ }\PY{k}{as}\PY{+w}{ }\PY{n+nn}{pd}
        \end{Verbatim}
    \end{tcolorbox}

    \begin{tcolorbox}[breakable, size=fbox, boxrule=1pt, pad at break*=1mm,colback=cellbackground, colframe=cellborder]
        \prompt{In}{incolor}{2}{\boxspacing}
        \begin{Verbatim}[commandchars=\\\{\}]
\PY{n}{dataset} \PY{o}{=} \PY{n}{pd}\PY{o}{.}\PY{n}{read\PYZus{}csv}\PY{p}{(}\PY{l+s+s2}{\PYZdq{}}\PY{l+s+s2}{Data.csv}\PY{l+s+s2}{\PYZdq{}}\PY{p}{)}
\PY{n}{dataset}
        \end{Verbatim}
    \end{tcolorbox}

    \begin{tcolorbox}[breakable, size=fbox, boxrule=.5pt, pad at break*=1mm, opacityfill=0]
        \prompt{Out}{outcolor}{2}{\boxspacing}
        \begin{Verbatim}[commandchars=\\\{\}]
   Country   Age   Salary Purchased
0   France  44.0  72000.0        No
1    Spain  27.0  48000.0       Yes
2  Germany  30.0  54000.0        No
3    Spain  38.0  61000.0        No
4  Germany  40.0      NaN       Yes
5   France  35.0  58000.0       Yes
6    Spain   NaN  52000.0        No
7   France  48.0  79000.0       Yes
8  Germany  50.0  83000.0        No
9   France  37.0  67000.0       Yes
        \end{Verbatim}
    \end{tcolorbox}

    \begin{tcolorbox}[breakable, size=fbox, boxrule=1pt, pad at break*=1mm,colback=cellbackground, colframe=cellborder]
        \prompt{In}{incolor}{3}{\boxspacing}
        \begin{Verbatim}[commandchars=\\\{\}]
\PY{c+c1}{\PYZsh{} numpy}
\PY{n}{X\PYZus{}raw} \PY{o}{=} \PY{n}{dataset}\PY{o}{.}\PY{n}{iloc}\PY{p}{[}\PY{p}{:}\PY{p}{,} \PY{p}{:}\PY{o}{\PYZhy{}}\PY{l+m+mi}{1}\PY{p}{]}\PY{o}{.}\PY{n}{values}
\PY{n}{y\PYZus{}raw} \PY{o}{=} \PY{n}{dataset}\PY{o}{.}\PY{n}{iloc}\PY{p}{[}\PY{p}{:}\PY{p}{,} \PY{o}{\PYZhy{}}\PY{l+m+mi}{1}\PY{p}{]}\PY{o}{.}\PY{n}{values}
        \end{Verbatim}
    \end{tcolorbox}

    \begin{tcolorbox}[breakable, size=fbox, boxrule=1pt, pad at break*=1mm,colback=cellbackground, colframe=cellborder]
        \prompt{In}{incolor}{4}{\boxspacing}
        \begin{Verbatim}[commandchars=\\\{\}]
\PY{n+nb}{print}\PY{p}{(}\PY{n}{X\PYZus{}raw}\PY{p}{)}
\PY{n+nb}{print}\PY{p}{(}\PY{n}{y\PYZus{}raw}\PY{p}{)}
        \end{Verbatim}
    \end{tcolorbox}

    \begin{Verbatim}[commandchars=\\\{\}]
[['France' 44.0 72000.0]
 ['Spain' 27.0 48000.0]
 ['Germany' 30.0 54000.0]
 ['Spain' 38.0 61000.0]
 ['Germany' 40.0 nan]
 ['France' 35.0 58000.0]
 ['Spain' nan 52000.0]
 ['France' 48.0 79000.0]
 ['Germany' 50.0 83000.0]
 ['France' 37.0 67000.0]]
<StringArray>
['No', 'Yes', 'No', 'No', 'Yes', 'Yes', 'No', 'Yes', 'No', 'Yes']
Length: 10, dtype: str
    \end{Verbatim}

    \begin{tcolorbox}[breakable, size=fbox, boxrule=1pt, pad at break*=1mm,colback=cellbackground, colframe=cellborder]
        \prompt{In}{incolor}{5}{\boxspacing}
        \begin{Verbatim}[commandchars=\\\{\}]
\PY{c+c1}{\PYZsh{} pandas}
\PY{n}{dfX\PYZus{}raw} \PY{o}{=} \PY{n}{dataset}\PY{o}{.}\PY{n}{iloc}\PY{p}{[}\PY{p}{:}\PY{p}{,} \PY{p}{:}\PY{o}{\PYZhy{}}\PY{l+m+mi}{1}\PY{p}{]}
\PY{n}{seriesy\PYZus{}rawdfX} \PY{o}{=} \PY{n}{dataset}\PY{o}{.}\PY{n}{iloc}\PY{p}{[}\PY{p}{:}\PY{p}{,} \PY{o}{\PYZhy{}}\PY{l+m+mi}{1}\PY{p}{]}
        \end{Verbatim}
    \end{tcolorbox}

    \begin{tcolorbox}[breakable, size=fbox, boxrule=1pt, pad at break*=1mm,colback=cellbackground, colframe=cellborder]
        \prompt{In}{incolor}{6}{\boxspacing}
        \begin{Verbatim}[commandchars=\\\{\}]
\PY{n+nb}{print}\PY{p}{(}\PY{n}{dfX\PYZus{}raw}\PY{p}{)}
\PY{n+nb}{print}\PY{p}{(}\PY{n}{seriesy\PYZus{}rawdfX}\PY{p}{)}
        \end{Verbatim}
    \end{tcolorbox}

    \begin{Verbatim}[commandchars=\\\{\}]
   Country   Age   Salary
0   France  44.0  72000.0
1    Spain  27.0  48000.0
2  Germany  30.0  54000.0
3    Spain  38.0  61000.0
4  Germany  40.0      NaN
5   France  35.0  58000.0
6    Spain   NaN  52000.0
7   France  48.0  79000.0
8  Germany  50.0  83000.0
9   France  37.0  67000.0
0     No
1    Yes
2     No
3     No
4    Yes
5    Yes
6     No
7    Yes
8     No
9    Yes
Name: Purchased, dtype: str
    \end{Verbatim}

    \begin{tcolorbox}[breakable, size=fbox, boxrule=1pt, pad at break*=1mm,colback=cellbackground, colframe=cellborder]
        \prompt{In}{incolor}{7}{\boxspacing}
        \begin{Verbatim}[commandchars=\\\{\}]
\PY{c+c1}{\PYZsh{} copy the data to another variable}
\PY{n}{dfX} \PY{o}{=} \PY{n}{dfX\PYZus{}raw}\PY{o}{.}\PY{n}{copy}\PY{p}{(}\PY{p}{)}
\PY{n}{seriesy} \PY{o}{=} \PY{n}{seriesy\PYZus{}rawdfX}\PY{o}{.}\PY{n}{copy}\PY{p}{(}\PY{p}{)}
\PY{n}{X} \PY{o}{=} \PY{n}{X\PYZus{}raw}\PY{o}{.}\PY{n}{copy}\PY{p}{(}\PY{p}{)}
\PY{n}{y} \PY{o}{=} \PY{n}{y\PYZus{}raw}\PY{o}{.}\PY{n}{copy}\PY{p}{(}\PY{p}{)}
        \end{Verbatim}
    \end{tcolorbox}

    \begin{tcolorbox}[breakable, size=fbox, boxrule=1pt, pad at break*=1mm,colback=cellbackground, colframe=cellborder]
        \prompt{In}{incolor}{8}{\boxspacing}
        \begin{Verbatim}[commandchars=\\\{\}]
\PY{n+nb}{print}\PY{p}{(}\PY{n}{X}\PY{p}{)}
        \end{Verbatim}
    \end{tcolorbox}

    \begin{Verbatim}[commandchars=\\\{\}]
[['France' 44.0 72000.0]
 ['Spain' 27.0 48000.0]
 ['Germany' 30.0 54000.0]
 ['Spain' 38.0 61000.0]
 ['Germany' 40.0 nan]
 ['France' 35.0 58000.0]
 ['Spain' nan 52000.0]
 ['France' 48.0 79000.0]
 ['Germany' 50.0 83000.0]
 ['France' 37.0 67000.0]]
    \end{Verbatim}

    \begin{tcolorbox}[breakable, size=fbox, boxrule=1pt, pad at break*=1mm,colback=cellbackground, colframe=cellborder]
        \prompt{In}{incolor}{9}{\boxspacing}
        \begin{Verbatim}[commandchars=\\\{\}]
\PY{n+nb}{print}\PY{p}{(}\PY{n+nb}{type}\PY{p}{(}\PY{n}{dfX}\PY{p}{)}\PY{p}{)}
\PY{n+nb}{print}\PY{p}{(}\PY{n+nb}{type}\PY{p}{(}\PY{n}{X}\PY{p}{)}\PY{p}{)}
        \end{Verbatim}
    \end{tcolorbox}

    \begin{Verbatim}[commandchars=\\\{\}]
<class 'pandas.DataFrame'>
<class 'numpy.ndarray'>
    \end{Verbatim}

    \begin{tcolorbox}[breakable, size=fbox, boxrule=1pt, pad at break*=1mm,colback=cellbackground, colframe=cellborder]
        \prompt{In}{incolor}{10}{\boxspacing}
        \begin{Verbatim}[commandchars=\\\{\}]
\PY{n+nb}{print}\PY{p}{(}\PY{n}{y}\PY{p}{)}
        \end{Verbatim}
    \end{tcolorbox}

    \begin{Verbatim}[commandchars=\\\{\}]
<StringArray>
['No', 'Yes', 'No', 'No', 'Yes', 'Yes', 'No', 'Yes', 'No', 'Yes']
Length: 10, dtype: str
    \end{Verbatim}

    \begin{tcolorbox}[breakable, size=fbox, boxrule=1pt, pad at break*=1mm,colback=cellbackground, colframe=cellborder]
        \prompt{In}{incolor}{11}{\boxspacing}
        \begin{Verbatim}[commandchars=\\\{\}]
\PY{n+nb}{print}\PY{p}{(}\PY{n+nb}{type}\PY{p}{(}\PY{n}{seriesy}\PY{p}{)}\PY{p}{)}
\PY{n+nb}{print}\PY{p}{(}\PY{n+nb}{type}\PY{p}{(}\PY{n}{y}\PY{p}{)}\PY{p}{)}
        \end{Verbatim}
    \end{tcolorbox}

    \begin{Verbatim}[commandchars=\\\{\}]
<class 'pandas.Series'>
<class 'pandas.arrays.StringArray'>
    \end{Verbatim}

    \hypertarget{taking-care-of-missing-data}{%
        \subsection{Taking care of missing
        data}\label{taking-care-of-missing-data}}

    \begin{tcolorbox}[breakable, size=fbox, boxrule=1pt, pad at break*=1mm,colback=cellbackground, colframe=cellborder]
        \prompt{In}{incolor}{12}{\boxspacing}
        \begin{Verbatim}[commandchars=\\\{\}]
\PY{n+nb}{print}\PY{p}{(}\PY{n}{dfX}\PY{o}{.}\PY{n}{info}\PY{p}{(}\PY{p}{)}\PY{p}{)}
        \end{Verbatim}
    \end{tcolorbox}

    \begin{Verbatim}[commandchars=\\\{\}]
<class 'pandas.DataFrame'>
RangeIndex: 10 entries, 0 to 9
Data columns (total 3 columns):
 \#   Column   Non-Null Count  Dtype
---  ------   --------------  -----
 0   Country  10 non-null     str
 1   Age      9 non-null      float64
 2   Salary   9 non-null      float64
dtypes: float64(2), str(1)
memory usage: 372.0 bytes
None
    \end{Verbatim}

    \begin{tcolorbox}[breakable, size=fbox, boxrule=1pt, pad at break*=1mm,colback=cellbackground, colframe=cellborder]
        \prompt{In}{incolor}{13}{\boxspacing}
        \begin{Verbatim}[commandchars=\\\{\}]
\PY{n+nb}{print}\PY{p}{(}\PY{n}{X}\PY{o}{.}\PY{n}{info}\PY{p}{(}\PY{p}{)}\PY{p}{)}
        \end{Verbatim}
    \end{tcolorbox}

    \begin{Verbatim}[commandchars=\\\{\}, frame=single, framerule=2mm, rulecolor=\color{outerrorbackground}]
\textcolor{ansi-red}{---------------------------------------------------------------------------}
\textcolor{ansi-red}{AttributeError}                            Traceback (most recent call last)
\textcolor{ansi-cyan}{Cell}\textcolor{ansi-cyan}{ }\textcolor{ansi-green}{In[13]}\textcolor{ansi-green}{, line 1}
\textcolor{ansi-green}{----> }\textcolor{ansi-green}{1} print(X.info())

\textcolor{ansi-red}{AttributeError}: 'numpy.ndarray' object has no attribute 'info'
    \end{Verbatim}

    \begin{tcolorbox}[breakable, size=fbox, boxrule=1pt, pad at break*=1mm,colback=cellbackground, colframe=cellborder]
        \prompt{In}{incolor}{14}{\boxspacing}
        \begin{Verbatim}[commandchars=\\\{\}]
\PY{c+c1}{\PYZsh{} change numpy array to dataframe}
\PY{n}{X} \PY{o}{=} \PY{n}{pd}\PY{o}{.}\PY{n}{DataFrame}\PY{p}{(}\PY{n}{X}\PY{p}{)}
\PY{n}{X}\PY{o}{.}\PY{n}{info}\PY{p}{(}\PY{p}{)}
        \end{Verbatim}
    \end{tcolorbox}

    \begin{Verbatim}[commandchars=\\\{\}]
<class 'pandas.DataFrame'>
RangeIndex: 10 entries, 0 to 9
Data columns (total 3 columns):
 \#   Column  Non-Null Count  Dtype
---  ------  --------------  -----
 0   0       10 non-null     str
 1   1       9 non-null      object
 2   2       9 non-null      object
dtypes: object(2), str(1)
memory usage: 372.0+ bytes
    \end{Verbatim}

    \begin{tcolorbox}[breakable, size=fbox, boxrule=1pt, pad at break*=1mm,colback=cellbackground, colframe=cellborder]
        \prompt{In}{incolor}{15}{\boxspacing}
        \begin{Verbatim}[commandchars=\\\{\}]
\PY{c+c1}{\PYZsh{} Then change to numpy}
\PY{n}{X} \PY{o}{=} \PY{n}{X}\PY{o}{.}\PY{n}{values}
\PY{n}{X}
        \end{Verbatim}
    \end{tcolorbox}

    \begin{tcolorbox}[breakable, size=fbox, boxrule=.5pt, pad at break*=1mm, opacityfill=0]
        \prompt{Out}{outcolor}{15}{\boxspacing}
        \begin{Verbatim}[commandchars=\\\{\}]
array([['France', 44.0, 72000.0],
       ['Spain', 27.0, 48000.0],
       ['Germany', 30.0, 54000.0],
       ['Spain', 38.0, 61000.0],
       ['Germany', 40.0, nan],
       ['France', 35.0, 58000.0],
       ['Spain', nan, 52000.0],
       ['France', 48.0, 79000.0],
       ['Germany', 50.0, 83000.0],
       ['France', 37.0, 67000.0]], dtype=object)
        \end{Verbatim}
    \end{tcolorbox}

    \begin{tcolorbox}[breakable, size=fbox, boxrule=1pt, pad at break*=1mm,colback=cellbackground, colframe=cellborder]
        \prompt{In}{incolor}{16}{\boxspacing}
        \begin{Verbatim}[commandchars=\\\{\}]
\PY{n+nb}{print}\PY{p}{(}\PY{n}{dfX}\PY{o}{.}\PY{n}{isnull}\PY{p}{(}\PY{p}{)}\PY{o}{.}\PY{n}{sum}\PY{p}{(}\PY{p}{)}\PY{p}{)}
        \end{Verbatim}
    \end{tcolorbox}

    \begin{Verbatim}[commandchars=\\\{\}]
Country    0
Age        1
Salary     1
dtype: int64
    \end{Verbatim}

    \begin{tcolorbox}[breakable, size=fbox, boxrule=1pt, pad at break*=1mm,colback=cellbackground, colframe=cellborder]
        \prompt{In}{incolor}{17}{\boxspacing}
        \begin{Verbatim}[commandchars=\\\{\}]
\PY{n+nb}{print}\PY{p}{(}\PY{n}{X}\PY{o}{.}\PY{n}{isnull}\PY{p}{(}\PY{p}{)}\PY{o}{.}\PY{n}{sum}\PY{p}{(}\PY{p}{)}\PY{p}{)}
        \end{Verbatim}
    \end{tcolorbox}

    \begin{Verbatim}[commandchars=\\\{\}, frame=single, framerule=2mm, rulecolor=\color{outerrorbackground}]
\textcolor{ansi-red}{---------------------------------------------------------------------------}
\textcolor{ansi-red}{AttributeError}                            Traceback (most recent call last)
\textcolor{ansi-cyan}{Cell}\textcolor{ansi-cyan}{ }\textcolor{ansi-green}{In[17]}\textcolor{ansi-green}{, line 1}
\textcolor{ansi-green}{----> }\textcolor{ansi-green}{1} print(X.isnull().sum())

\textcolor{ansi-red}{AttributeError}: 'numpy.ndarray' object has no attribute 'isnull'
    \end{Verbatim}

    \begin{tcolorbox}[breakable, size=fbox, boxrule=1pt, pad at break*=1mm,colback=cellbackground, colframe=cellborder]
        \prompt{In}{incolor}{18}{\boxspacing}
        \begin{Verbatim}[commandchars=\\\{\}]
\PY{n}{dfX\PYZus{}raw}
        \end{Verbatim}
    \end{tcolorbox}

    \begin{tcolorbox}[breakable, size=fbox, boxrule=.5pt, pad at break*=1mm, opacityfill=0]
        \prompt{Out}{outcolor}{18}{\boxspacing}
        \begin{Verbatim}[commandchars=\\\{\}]
   Country   Age   Salary
0   France  44.0  72000.0
1    Spain  27.0  48000.0
2  Germany  30.0  54000.0
3    Spain  38.0  61000.0
4  Germany  40.0      NaN
5   France  35.0  58000.0
6    Spain   NaN  52000.0
7   France  48.0  79000.0
8  Germany  50.0  83000.0
9   France  37.0  67000.0
        \end{Verbatim}
    \end{tcolorbox}

    \begin{tcolorbox}[breakable, size=fbox, boxrule=1pt, pad at break*=1mm,colback=cellbackground, colframe=cellborder]
        \prompt{In}{incolor}{19}{\boxspacing}
        \begin{Verbatim}[commandchars=\\\{\}]
\PY{n}{dfX}
        \end{Verbatim}
    \end{tcolorbox}

    \begin{tcolorbox}[breakable, size=fbox, boxrule=.5pt, pad at break*=1mm, opacityfill=0]
        \prompt{Out}{outcolor}{19}{\boxspacing}
        \begin{Verbatim}[commandchars=\\\{\}]
   Country   Age   Salary
0   France  44.0  72000.0
1    Spain  27.0  48000.0
2  Germany  30.0  54000.0
3    Spain  38.0  61000.0
4  Germany  40.0      NaN
5   France  35.0  58000.0
6    Spain   NaN  52000.0
7   France  48.0  79000.0
8  Germany  50.0  83000.0
9   France  37.0  67000.0
        \end{Verbatim}
    \end{tcolorbox}

    \begin{tcolorbox}[breakable, size=fbox, boxrule=1pt, pad at break*=1mm,colback=cellbackground, colframe=cellborder]
        \prompt{In}{incolor}{20}{\boxspacing}
        \begin{Verbatim}[commandchars=\\\{\}]
\PY{n}{dfX} \PY{o}{=} \PY{n}{dfX}\PY{o}{.}\PY{n}{fillna}\PY{p}{(}\PY{l+m+mi}{0}\PY{p}{)}
\PY{n}{dfX}
        \end{Verbatim}
    \end{tcolorbox}

    \begin{tcolorbox}[breakable, size=fbox, boxrule=.5pt, pad at break*=1mm, opacityfill=0]
        \prompt{Out}{outcolor}{20}{\boxspacing}
        \begin{Verbatim}[commandchars=\\\{\}]
   Country   Age   Salary
0   France  44.0  72000.0
1    Spain  27.0  48000.0
2  Germany  30.0  54000.0
3    Spain  38.0  61000.0
4  Germany  40.0      0.0
5   France  35.0  58000.0
6    Spain   0.0  52000.0
7   France  48.0  79000.0
8  Germany  50.0  83000.0
9   France  37.0  67000.0
        \end{Verbatim}
    \end{tcolorbox}

    \begin{tcolorbox}[breakable, size=fbox, boxrule=1pt, pad at break*=1mm,colback=cellbackground, colframe=cellborder]
        \prompt{In}{incolor}{21}{\boxspacing}
        \begin{Verbatim}[commandchars=\\\{\}]
\PY{c+c1}{\PYZsh{} replacing 0 with NaN value}

\PY{n}{dfX} \PY{o}{=} \PY{n}{dfX}\PY{o}{.}\PY{n}{replace}\PY{p}{(}\PY{l+m+mi}{0}\PY{p}{,} \PY{n}{np}\PY{o}{.}\PY{n}{nan}\PY{p}{)}
\PY{n}{dfX}
        \end{Verbatim}
    \end{tcolorbox}

    \begin{tcolorbox}[breakable, size=fbox, boxrule=.5pt, pad at break*=1mm, opacityfill=0]
        \prompt{Out}{outcolor}{21}{\boxspacing}
        \begin{Verbatim}[commandchars=\\\{\}]
   Country   Age   Salary
0   France  44.0  72000.0
1    Spain  27.0  48000.0
2  Germany  30.0  54000.0
3    Spain  38.0  61000.0
4  Germany  40.0      NaN
5   France  35.0  58000.0
6    Spain   NaN  52000.0
7   France  48.0  79000.0
8  Germany  50.0  83000.0
9   France  37.0  67000.0
        \end{Verbatim}
    \end{tcolorbox}

    \begin{tcolorbox}[breakable, size=fbox, boxrule=1pt, pad at break*=1mm,colback=cellbackground, colframe=cellborder]
        \prompt{In}{incolor}{22}{\boxspacing}
        \begin{Verbatim}[commandchars=\\\{\}]
\PY{n}{index\PYZus{}no} \PY{o}{=} \PY{n}{dfX}\PY{o}{.}\PY{n}{columns}\PY{o}{.}\PY{n}{get\PYZus{}loc}\PY{p}{(}\PY{l+s+s2}{\PYZdq{}}\PY{l+s+s2}{Age}\PY{l+s+s2}{\PYZdq{}}\PY{p}{)}
\PY{n}{index\PYZus{}no}
        \end{Verbatim}
    \end{tcolorbox}

    \begin{tcolorbox}[breakable, size=fbox, boxrule=.5pt, pad at break*=1mm, opacityfill=0]
        \prompt{Out}{outcolor}{22}{\boxspacing}
        \begin{Verbatim}[commandchars=\\\{\}]
1
        \end{Verbatim}
    \end{tcolorbox}

    \begin{tcolorbox}[breakable, size=fbox, boxrule=1pt, pad at break*=1mm,colback=cellbackground, colframe=cellborder]
        \prompt{In}{incolor}{23}{\boxspacing}
        \begin{Verbatim}[commandchars=\\\{\}]
\PY{n}{dfX}\PY{o}{.}\PY{n}{iloc}\PY{p}{[}\PY{p}{:}\PY{p}{,} \PY{l+m+mi}{1}\PY{p}{]} \PY{o}{=} \PY{n}{dfX}\PY{o}{.}\PY{n}{iloc}\PY{p}{[}\PY{p}{:}\PY{p}{,} \PY{l+m+mi}{1}\PY{p}{]}\PY{o}{.}\PY{n}{fillna}\PY{p}{(}\PY{l+m+mi}{5}\PY{p}{)}
\PY{n}{dfX}
        \end{Verbatim}
    \end{tcolorbox}

    \begin{tcolorbox}[breakable, size=fbox, boxrule=.5pt, pad at break*=1mm, opacityfill=0]
        \prompt{Out}{outcolor}{23}{\boxspacing}
        \begin{Verbatim}[commandchars=\\\{\}]
   Country   Age   Salary
0   France  44.0  72000.0
1    Spain  27.0  48000.0
2  Germany  30.0  54000.0
3    Spain  38.0  61000.0
4  Germany  40.0      NaN
5   France  35.0  58000.0
6    Spain   5.0  52000.0
7   France  48.0  79000.0
8  Germany  50.0  83000.0
9   France  37.0  67000.0
        \end{Verbatim}
    \end{tcolorbox}

    \begin{tcolorbox}[breakable, size=fbox, boxrule=1pt, pad at break*=1mm,colback=cellbackground, colframe=cellborder]
        \prompt{In}{incolor}{24}{\boxspacing}
        \begin{Verbatim}[commandchars=\\\{\}]
\PY{c+c1}{\PYZsh{} getting the old data}
\PY{n}{dfX} \PY{o}{=} \PY{n}{dfX\PYZus{}raw}\PY{o}{.}\PY{n}{copy}\PY{p}{(}\PY{p}{)}
\PY{n}{dfX}
        \end{Verbatim}
    \end{tcolorbox}

    \begin{tcolorbox}[breakable, size=fbox, boxrule=.5pt, pad at break*=1mm, opacityfill=0]
        \prompt{Out}{outcolor}{24}{\boxspacing}
        \begin{Verbatim}[commandchars=\\\{\}]
   Country   Age   Salary
0   France  44.0  72000.0
1    Spain  27.0  48000.0
2  Germany  30.0  54000.0
3    Spain  38.0  61000.0
4  Germany  40.0      NaN
5   France  35.0  58000.0
6    Spain   NaN  52000.0
7   France  48.0  79000.0
8  Germany  50.0  83000.0
9   France  37.0  67000.0
        \end{Verbatim}
    \end{tcolorbox}

    \hypertarget{missing-value-imputation-in-dataframe}{%
        \subsection{Missing value imputation in
        dataframe}\label{missing-value-imputation-in-dataframe}}

    \hypertarget{replace-with-mean-values}{%
        \subsubsection{Replace with mean
        values}\label{replace-with-mean-values}}

    \begin{tcolorbox}[breakable, size=fbox, boxrule=1pt, pad at break*=1mm,colback=cellbackground, colframe=cellborder]
        \prompt{In}{incolor}{25}{\boxspacing}
        \begin{Verbatim}[commandchars=\\\{\}]
\PY{n}{dfX\PYZus{}raw}
        \end{Verbatim}
    \end{tcolorbox}

    \begin{tcolorbox}[breakable, size=fbox, boxrule=.5pt, pad at break*=1mm, opacityfill=0]
        \prompt{Out}{outcolor}{25}{\boxspacing}
        \begin{Verbatim}[commandchars=\\\{\}]
   Country   Age   Salary
0   France  44.0  72000.0
1    Spain  27.0  48000.0
2  Germany  30.0  54000.0
3    Spain  38.0  61000.0
4  Germany  40.0      NaN
5   France  35.0  58000.0
6    Spain   NaN  52000.0
7   France  48.0  79000.0
8  Germany  50.0  83000.0
9   France  37.0  67000.0
        \end{Verbatim}
    \end{tcolorbox}

    \begin{tcolorbox}[breakable, size=fbox, boxrule=1pt, pad at break*=1mm,colback=cellbackground, colframe=cellborder]
        \prompt{In}{incolor}{26}{\boxspacing}
        \begin{Verbatim}[commandchars=\\\{\}]
\PY{c+c1}{\PYZsh{} datafram}
\PY{n}{dfX} \PY{o}{=} \PY{n}{dfX\PYZus{}raw}\PY{o}{.}\PY{n}{copy}\PY{p}{(}\PY{p}{)}

\PY{k+kn}{from}\PY{+w}{ }\PY{n+nn}{sklearn}\PY{n+nn}{.}\PY{n+nn}{impute}\PY{+w}{ }\PY{k+kn}{import} \PY{n}{SimpleImputer}    \PY{c+c1}{\PYZsh{} SimpleImputer is a class ie usedd to handle the misssing values with simple strategies}
\PY{n}{imputer} \PY{o}{=} \PY{n}{SimpleImputer}\PY{p}{(}\PY{n}{missing\PYZus{}values}\PY{o}{=}\PY{n}{np}\PY{o}{.}\PY{n}{nan}\PY{p}{,} \PY{n}{strategy}\PY{o}{=}\PY{l+s+s2}{\PYZdq{}}\PY{l+s+s2}{mean}\PY{l+s+s2}{\PYZdq{}}\PY{p}{)}

\PY{c+c1}{\PYZsh{} imputer.fit(dfX.iloc[:, 1:3])}
\PY{c+c1}{\PYZsh{} dfX.iloc[:, 1:3] = imputer.transform(dfX.iloc[:, 1:3])}
\PY{c+c1}{\PYZsh{} dfX.iloc[;, 1:3] = imputer.fit(dfX.iloc[:, 1:3]).transform(dfX.iloc[:, 1:3])}

\PY{n}{dfX}\PY{o}{.}\PY{n}{iloc}\PY{p}{[}\PY{p}{:}\PY{p}{,} \PY{l+m+mi}{1}\PY{p}{:}\PY{l+m+mi}{3}\PY{p}{]} \PY{o}{=} \PY{n}{imputer}\PY{o}{.}\PY{n}{fit\PYZus{}transform}\PY{p}{(}\PY{n}{dfX}\PY{o}{.}\PY{n}{iloc}\PY{p}{[}\PY{p}{:}\PY{p}{,} \PY{l+m+mi}{1}\PY{p}{:}\PY{l+m+mi}{3}\PY{p}{]}\PY{p}{)}
\PY{n}{dfX}
        \end{Verbatim}
    \end{tcolorbox}

    \begin{tcolorbox}[breakable, size=fbox, boxrule=.5pt, pad at break*=1mm, opacityfill=0]
        \prompt{Out}{outcolor}{26}{\boxspacing}
        \begin{Verbatim}[commandchars=\\\{\}]
   Country        Age        Salary
0   France  44.000000  72000.000000
1    Spain  27.000000  48000.000000
2  Germany  30.000000  54000.000000
3    Spain  38.000000  61000.000000
4  Germany  40.000000  63777.777778
5   France  35.000000  58000.000000
6    Spain  38.777778  52000.000000
7   France  48.000000  79000.000000
8  Germany  50.000000  83000.000000
9   France  37.000000  67000.000000
        \end{Verbatim}
    \end{tcolorbox}

    \begin{tcolorbox}[breakable, size=fbox, boxrule=1pt, pad at break*=1mm,colback=cellbackground, colframe=cellborder]
        \prompt{In}{incolor}{27}{\boxspacing}
        \begin{Verbatim}[commandchars=\\\{\}]
\PY{c+c1}{\PYZsh{} datafram}
\PY{n}{dfX} \PY{o}{=} \PY{n}{dfX\PYZus{}raw}\PY{o}{.}\PY{n}{copy}\PY{p}{(}\PY{p}{)}

\PY{k+kn}{from}\PY{+w}{ }\PY{n+nn}{sklearn}\PY{n+nn}{.}\PY{n+nn}{impute}\PY{+w}{ }\PY{k+kn}{import} \PY{n}{SimpleImputer}    \PY{c+c1}{\PYZsh{} SimpleImputer is a class ie usedd to handle the misssing values with simple strategies}
\PY{n}{imputer} \PY{o}{=} \PY{n}{SimpleImputer}\PY{p}{(}\PY{n}{missing\PYZus{}values}\PY{o}{=}\PY{n}{np}\PY{o}{.}\PY{n}{nan}\PY{p}{,} \PY{n}{strategy}\PY{o}{=}\PY{l+s+s2}{\PYZdq{}}\PY{l+s+s2}{median}\PY{l+s+s2}{\PYZdq{}}\PY{p}{)}

\PY{c+c1}{\PYZsh{} imputer.fit(dfX.iloc[:, 1:3])}
\PY{c+c1}{\PYZsh{} dfX.iloc[:, 1:3] = imputer.transform(dfX.iloc[:, 1:3])}
\PY{c+c1}{\PYZsh{} dfX.iloc[;, 1:3] = imputer.fit(dfX.iloc[:, 1:3]).transform(dfX.iloc[:, 1:3])}

\PY{n}{dfX}\PY{o}{.}\PY{n}{iloc}\PY{p}{[}\PY{p}{:}\PY{p}{,} \PY{l+m+mi}{1}\PY{p}{:}\PY{l+m+mi}{3}\PY{p}{]} \PY{o}{=} \PY{n}{imputer}\PY{o}{.}\PY{n}{fit\PYZus{}transform}\PY{p}{(}\PY{n}{dfX}\PY{o}{.}\PY{n}{iloc}\PY{p}{[}\PY{p}{:}\PY{p}{,} \PY{l+m+mi}{1}\PY{p}{:}\PY{l+m+mi}{3}\PY{p}{]}\PY{p}{)}
\PY{n}{dfX}
        \end{Verbatim}
    \end{tcolorbox}

    \begin{tcolorbox}[breakable, size=fbox, boxrule=.5pt, pad at break*=1mm, opacityfill=0]
        \prompt{Out}{outcolor}{27}{\boxspacing}
        \begin{Verbatim}[commandchars=\\\{\}]
   Country   Age   Salary
0   France  44.0  72000.0
1    Spain  27.0  48000.0
2  Germany  30.0  54000.0
3    Spain  38.0  61000.0
4  Germany  40.0  61000.0
5   France  35.0  58000.0
6    Spain  38.0  52000.0
7   France  48.0  79000.0
8  Germany  50.0  83000.0
9   France  37.0  67000.0
        \end{Verbatim}
    \end{tcolorbox}

    \begin{tcolorbox}[breakable, size=fbox, boxrule=1pt, pad at break*=1mm,colback=cellbackground, colframe=cellborder]
        \prompt{In}{incolor}{28}{\boxspacing}
        \begin{Verbatim}[commandchars=\\\{\}]
\PY{c+c1}{\PYZsh{} getting the old data}
\PY{n}{dfX} \PY{o}{=} \PY{n}{dfX\PYZus{}raw}\PY{o}{.}\PY{n}{copy}\PY{p}{(}\PY{p}{)}
\PY{n}{dfX}
        \end{Verbatim}
    \end{tcolorbox}

    \begin{tcolorbox}[breakable, size=fbox, boxrule=.5pt, pad at break*=1mm, opacityfill=0]
        \prompt{Out}{outcolor}{28}{\boxspacing}
        \begin{Verbatim}[commandchars=\\\{\}]
   Country   Age   Salary
0   France  44.0  72000.0
1    Spain  27.0  48000.0
2  Germany  30.0  54000.0
3    Spain  38.0  61000.0
4  Germany  40.0      NaN
5   France  35.0  58000.0
6    Spain   NaN  52000.0
7   France  48.0  79000.0
8  Germany  50.0  83000.0
9   France  37.0  67000.0
        \end{Verbatim}
    \end{tcolorbox}

    \begin{tcolorbox}[breakable, size=fbox, boxrule=1pt, pad at break*=1mm,colback=cellbackground, colframe=cellborder]
        \prompt{In}{incolor}{29}{\boxspacing}
        \begin{Verbatim}[commandchars=\\\{\}]
\PY{c+c1}{\PYZsh{} dataframe}

\PY{k+kn}{from}\PY{+w}{ }\PY{n+nn}{sklearn}\PY{n+nn}{.}\PY{n+nn}{impute}\PY{+w}{ }\PY{k+kn}{import} \PY{n}{SimpleImputer}    \PY{c+c1}{\PYZsh{} SimpleImputer is a class ie usedd to handle the misssing values with simple strategies}
\PY{n}{imputer} \PY{o}{=} \PY{n}{SimpleImputer}\PY{p}{(}\PY{n}{missing\PYZus{}values}\PY{o}{=}\PY{n}{np}\PY{o}{.}\PY{n}{nan}\PY{p}{,} \PY{n}{strategy}\PY{o}{=}\PY{l+s+s2}{\PYZdq{}}\PY{l+s+s2}{constant}\PY{l+s+s2}{\PYZdq{}}\PY{p}{,} \PY{n}{fill\PYZus{}value}\PY{o}{=}\PY{l+m+mi}{5}\PY{p}{)}
\PY{n}{imputer}\PY{o}{.}\PY{n}{fit}\PY{p}{(}\PY{n}{dfX}\PY{o}{.}\PY{n}{iloc}\PY{p}{[}\PY{p}{:}\PY{p}{,} \PY{l+m+mi}{1}\PY{p}{:}\PY{l+m+mi}{3}\PY{p}{]}\PY{p}{)}
\PY{n}{dfX}\PY{o}{.}\PY{n}{iloc}\PY{p}{[}\PY{p}{:}\PY{p}{,} \PY{l+m+mi}{1}\PY{p}{:}\PY{l+m+mi}{3}\PY{p}{]} \PY{o}{=} \PY{n}{imputer}\PY{o}{.}\PY{n}{fit}\PY{p}{(}\PY{n}{dfX}\PY{o}{.}\PY{n}{iloc}\PY{p}{[}\PY{p}{:}\PY{p}{,} \PY{l+m+mi}{1}\PY{p}{:}\PY{l+m+mi}{3}\PY{p}{]}\PY{p}{)}\PY{o}{.}\PY{n}{transform}\PY{p}{(}\PY{n}{dfX}\PY{o}{.}\PY{n}{iloc}\PY{p}{[}\PY{p}{:}\PY{p}{,} \PY{l+m+mi}{1}\PY{p}{:}\PY{l+m+mi}{3}\PY{p}{]}\PY{p}{)}
\PY{n}{dfX}
        \end{Verbatim}
    \end{tcolorbox}

    \begin{tcolorbox}[breakable, size=fbox, boxrule=.5pt, pad at break*=1mm, opacityfill=0]
        \prompt{Out}{outcolor}{29}{\boxspacing}
        \begin{Verbatim}[commandchars=\\\{\}]
   Country   Age   Salary
0   France  44.0  72000.0
1    Spain  27.0  48000.0
2  Germany  30.0  54000.0
3    Spain  38.0  61000.0
4  Germany  40.0      5.0
5   France  35.0  58000.0
6    Spain   5.0  52000.0
7   France  48.0  79000.0
8  Germany  50.0  83000.0
9   France  37.0  67000.0
        \end{Verbatim}
    \end{tcolorbox}

    \hypertarget{missing-value-imputaion-with-numpy}{%
        \subsection{Missing value imputaion with
        numpy}\label{missing-value-imputaion-with-numpy}}

    \hypertarget{replace-with-mean-values}{%
        \subsubsection{Replace with mean
        values}\label{replace-with-mean-values}}

    \begin{tcolorbox}[breakable, size=fbox, boxrule=1pt, pad at break*=1mm,colback=cellbackground, colframe=cellborder]
        \prompt{In}{incolor}{30}{\boxspacing}
        \begin{Verbatim}[commandchars=\\\{\}]
\PY{n}{X\PYZus{}raw}
        \end{Verbatim}
    \end{tcolorbox}

    \begin{tcolorbox}[breakable, size=fbox, boxrule=.5pt, pad at break*=1mm, opacityfill=0]
        \prompt{Out}{outcolor}{30}{\boxspacing}
        \begin{Verbatim}[commandchars=\\\{\}]
array([['France', 44.0, 72000.0],
       ['Spain', 27.0, 48000.0],
       ['Germany', 30.0, 54000.0],
       ['Spain', 38.0, 61000.0],
       ['Germany', 40.0, nan],
       ['France', 35.0, 58000.0],
       ['Spain', nan, 52000.0],
       ['France', 48.0, 79000.0],
       ['Germany', 50.0, 83000.0],
       ['France', 37.0, 67000.0]], dtype=object)
        \end{Verbatim}
    \end{tcolorbox}

    \begin{tcolorbox}[breakable, size=fbox, boxrule=1pt, pad at break*=1mm,colback=cellbackground, colframe=cellborder]
        \prompt{In}{incolor}{31}{\boxspacing}
        \begin{Verbatim}[commandchars=\\\{\}]
\PY{n}{X} \PY{o}{=} \PY{n}{X\PYZus{}raw}\PY{o}{.}\PY{n}{copy}\PY{p}{(}\PY{p}{)}

\PY{k+kn}{from}\PY{+w}{ }\PY{n+nn}{sklearn}\PY{n+nn}{.}\PY{n+nn}{impute}\PY{+w}{ }\PY{k+kn}{import} \PY{n}{SimpleImputer}
\PY{n}{imputer} \PY{o}{=} \PY{n}{SimpleImputer}\PY{p}{(}\PY{n}{missing\PYZus{}values}\PY{o}{=}\PY{n}{np}\PY{o}{.}\PY{n}{nan}\PY{p}{,} \PY{n}{strategy}\PY{o}{=}\PY{l+s+s2}{\PYZdq{}}\PY{l+s+s2}{mean}\PY{l+s+s2}{\PYZdq{}}\PY{p}{)}
\PY{n}{imputer}\PY{o}{.}\PY{n}{fit}\PY{p}{(}\PY{n}{X}\PY{p}{[}\PY{p}{:}\PY{p}{,} \PY{l+m+mi}{1}\PY{p}{:}\PY{l+m+mi}{3}\PY{p}{]}\PY{p}{)}
\PY{n}{X}\PY{p}{[}\PY{p}{:}\PY{p}{,} \PY{l+m+mi}{1}\PY{p}{:}\PY{l+m+mi}{3}\PY{p}{]} \PY{o}{=} \PY{n}{imputer}\PY{o}{.}\PY{n}{fit\PYZus{}transform}\PY{p}{(}\PY{n}{X}\PY{p}{[}\PY{p}{:}\PY{p}{,} \PY{l+m+mi}{1}\PY{p}{:}\PY{l+m+mi}{3}\PY{p}{]}\PY{p}{)}
        \end{Verbatim}
    \end{tcolorbox}

    \begin{tcolorbox}[breakable, size=fbox, boxrule=1pt, pad at break*=1mm,colback=cellbackground, colframe=cellborder]
        \prompt{In}{incolor}{32}{\boxspacing}
        \begin{Verbatim}[commandchars=\\\{\}]
\PY{n+nb}{print}\PY{p}{(}\PY{n}{X}\PY{p}{)}
        \end{Verbatim}
    \end{tcolorbox}

    \begin{Verbatim}[commandchars=\\\{\}]
[['France' 44.0 72000.0]
 ['Spain' 27.0 48000.0]
 ['Germany' 30.0 54000.0]
 ['Spain' 38.0 61000.0]
 ['Germany' 40.0 63777.77777777778]
 ['France' 35.0 58000.0]
 ['Spain' 38.77777777777778 52000.0]
 ['France' 48.0 79000.0]
 ['Germany' 50.0 83000.0]
 ['France' 37.0 67000.0]]
    \end{Verbatim}

    \hypertarget{encoding-the-categorical-data}{%
        \subsection{Encoding the categorical
        data}\label{encoding-the-categorical-data}}

    \begin{tcolorbox}[breakable, size=fbox, boxrule=1pt, pad at break*=1mm,colback=cellbackground, colframe=cellborder]
        \prompt{In}{incolor}{33}{\boxspacing}
        \begin{Verbatim}[commandchars=\\\{\}]
\PY{n}{dfX}
        \end{Verbatim}
    \end{tcolorbox}

    \begin{tcolorbox}[breakable, size=fbox, boxrule=.5pt, pad at break*=1mm, opacityfill=0]
        \prompt{Out}{outcolor}{33}{\boxspacing}
        \begin{Verbatim}[commandchars=\\\{\}]
   Country   Age   Salary
0   France  44.0  72000.0
1    Spain  27.0  48000.0
2  Germany  30.0  54000.0
3    Spain  38.0  61000.0
4  Germany  40.0      5.0
5   France  35.0  58000.0
6    Spain   5.0  52000.0
7   France  48.0  79000.0
8  Germany  50.0  83000.0
9   France  37.0  67000.0
        \end{Verbatim}
    \end{tcolorbox}

    \begin{tcolorbox}[breakable, size=fbox, boxrule=1pt, pad at break*=1mm,colback=cellbackground, colframe=cellborder]
        \prompt{In}{incolor}{34}{\boxspacing}
        \begin{Verbatim}[commandchars=\\\{\}]
\PY{c+c1}{\PYZsh{} pandas}
\PY{n}{dfX} \PY{o}{=} \PY{n}{pd}\PY{o}{.}\PY{n}{get\PYZus{}dummies}\PY{p}{(}\PY{n}{dfX}\PY{p}{,} \PY{n}{columns}\PY{o}{=}\PY{p}{[}\PY{l+s+s2}{\PYZdq{}}\PY{l+s+s2}{Country}\PY{l+s+s2}{\PYZdq{}}\PY{p}{]}\PY{p}{)}  \PY{c+c1}{\PYZsh{} one\PYZhy{}hot encoding (for each category)}
\PY{n}{dfX}
        \end{Verbatim}
    \end{tcolorbox}

    \begin{tcolorbox}[breakable, size=fbox, boxrule=.5pt, pad at break*=1mm, opacityfill=0]
        \prompt{Out}{outcolor}{34}{\boxspacing}
        \begin{Verbatim}[commandchars=\\\{\}]
    Age   Salary  Country\_France  Country\_Germany  Country\_Spain
0  44.0  72000.0            True            False          False
1  27.0  48000.0           False            False           True
2  30.0  54000.0           False             True          False
3  38.0  61000.0           False            False           True
4  40.0      5.0           False             True          False
5  35.0  58000.0            True            False          False
6   5.0  52000.0           False            False           True
7  48.0  79000.0            True            False          False
8  50.0  83000.0           False             True          False
9  37.0  67000.0            True            False          False
        \end{Verbatim}
    \end{tcolorbox}

    \begin{tcolorbox}[breakable, size=fbox, boxrule=1pt, pad at break*=1mm,colback=cellbackground, colframe=cellborder]
        \prompt{In}{incolor}{35}{\boxspacing}
        \begin{Verbatim}[commandchars=\\\{\}]
\PY{k+kn}{from}\PY{+w}{ }\PY{n+nn}{sklearn}\PY{n+nn}{.}\PY{n+nn}{compose}\PY{+w}{ }\PY{k+kn}{import} \PY{n}{ColumnTransformer}
\PY{k+kn}{from}\PY{+w}{ }\PY{n+nn}{sklearn}\PY{n+nn}{.}\PY{n+nn}{preprocessing}\PY{+w}{ }\PY{k+kn}{import} \PY{n}{OneHotEncoder}
\PY{c+c1}{\PYZsh{} ct = ColumnTransformer(transformers=[(\PYZdq{}encode\PYZdq{}, \PYZdq{}drop\PYZdq{}, [0])], remainder=\PYZdq{}passthrough\PYZdq{})}
\PY{n}{ct} \PY{o}{=} \PY{n}{ColumnTransformer}\PY{p}{(}\PY{n}{transformers}\PY{o}{=}\PY{p}{[}\PY{p}{(}\PY{l+s+s2}{\PYZdq{}}\PY{l+s+s2}{encode}\PY{l+s+s2}{\PYZdq{}}\PY{p}{,} \PY{n}{OneHotEncoder}\PY{p}{(}\PY{p}{)}\PY{p}{,} \PY{p}{[}\PY{l+m+mi}{0}\PY{p}{]}\PY{p}{)}\PY{p}{]}\PY{p}{,} \PY{n}{remainder}\PY{o}{=}\PY{l+s+s2}{\PYZdq{}}\PY{l+s+s2}{passthrough}\PY{l+s+s2}{\PYZdq{}}\PY{p}{)}
\PY{c+c1}{\PYZsh{} ct = ColumnTransformer(transformers=[(\PYZdq{}encode\PYZdq{}, OneHotEncoder(), [0])])}
\PY{n}{X} \PY{o}{=} \PY{n}{np}\PY{o}{.}\PY{n}{array}\PY{p}{(}\PY{n}{ct}\PY{o}{.}\PY{n}{fit\PYZus{}transform}\PY{p}{(}\PY{n}{X}\PY{p}{)}\PY{p}{)}
        \end{Verbatim}
    \end{tcolorbox}

    \begin{tcolorbox}[breakable, size=fbox, boxrule=1pt, pad at break*=1mm,colback=cellbackground, colframe=cellborder]
        \prompt{In}{incolor}{36}{\boxspacing}
        \begin{Verbatim}[commandchars=\\\{\}]
\PY{n+nb}{print}\PY{p}{(}\PY{n}{X}\PY{p}{)}
        \end{Verbatim}
    \end{tcolorbox}

    \begin{Verbatim}[commandchars=\\\{\}]
[[1.0 0.0 0.0 44.0 72000.0]
 [0.0 0.0 1.0 27.0 48000.0]
 [0.0 1.0 0.0 30.0 54000.0]
 [0.0 0.0 1.0 38.0 61000.0]
 [0.0 1.0 0.0 40.0 63777.77777777778]
 [1.0 0.0 0.0 35.0 58000.0]
 [0.0 0.0 1.0 38.77777777777778 52000.0]
 [1.0 0.0 0.0 48.0 79000.0]
 [0.0 1.0 0.0 50.0 83000.0]
 [1.0 0.0 0.0 37.0 67000.0]]
    \end{Verbatim}

    \hypertarget{encoding-the-dependent-variable}{%
        \subsection{Encoding the Dependent
        Variable}\label{encoding-the-dependent-variable}}

    \begin{tcolorbox}[breakable, size=fbox, boxrule=1pt, pad at break*=1mm,colback=cellbackground, colframe=cellborder]
        \prompt{In}{incolor}{37}{\boxspacing}
        \begin{Verbatim}[commandchars=\\\{\}]
\PY{k+kn}{from}\PY{+w}{ }\PY{n+nn}{sklearn}\PY{n+nn}{.}\PY{n+nn}{preprocessing}\PY{+w}{ }\PY{k+kn}{import} \PY{n}{LabelEncoder}  \PY{c+c1}{\PYZsh{} used to convert the categorical target labels into numerical labels}
\PY{n}{le} \PY{o}{=} \PY{n}{LabelEncoder}\PY{p}{(}\PY{p}{)}
\PY{n}{y} \PY{o}{=} \PY{n}{le}\PY{o}{.}\PY{n}{fit\PYZus{}transform}\PY{p}{(}\PY{n}{y}\PY{p}{)}
        \end{Verbatim}
    \end{tcolorbox}

    \begin{tcolorbox}[breakable, size=fbox, boxrule=1pt, pad at break*=1mm,colback=cellbackground, colframe=cellborder]
        \prompt{In}{incolor}{38}{\boxspacing}
        \begin{Verbatim}[commandchars=\\\{\}]
\PY{n+nb}{print}\PY{p}{(}\PY{n}{y}\PY{p}{)}
        \end{Verbatim}
    \end{tcolorbox}

    \begin{Verbatim}[commandchars=\\\{\}]
[0 1 0 0 1 1 0 1 0 1]
    \end{Verbatim}

    \begin{tcolorbox}[breakable, size=fbox, boxrule=1pt, pad at break*=1mm,colback=cellbackground, colframe=cellborder]
        \prompt{In}{incolor}{39}{\boxspacing}
        \begin{Verbatim}[commandchars=\\\{\}]
\PY{n}{le}\PY{o}{.}\PY{n}{fit}\PY{p}{(}\PY{p}{[}\PY{l+s+s2}{\PYZdq{}}\PY{l+s+s2}{india}\PY{l+s+s2}{\PYZdq{}}\PY{p}{,} \PY{l+s+s2}{\PYZdq{}}\PY{l+s+s2}{paris}\PY{l+s+s2}{\PYZdq{}}\PY{p}{,} \PY{l+s+s2}{\PYZdq{}}\PY{l+s+s2}{paris}\PY{l+s+s2}{\PYZdq{}}\PY{p}{,} \PY{l+s+s2}{\PYZdq{}}\PY{l+s+s2}{tokyo}\PY{l+s+s2}{\PYZdq{}}\PY{p}{,} \PY{l+s+s2}{\PYZdq{}}\PY{l+s+s2}{amsterdam}\PY{l+s+s2}{\PYZdq{}}\PY{p}{]}\PY{p}{)}
\PY{n+nb}{print}\PY{p}{(}\PY{n}{le}\PY{o}{.}\PY{n}{transform}\PY{p}{(}\PY{p}{[}\PY{l+s+s2}{\PYZdq{}}\PY{l+s+s2}{india}\PY{l+s+s2}{\PYZdq{}}\PY{p}{,} \PY{l+s+s2}{\PYZdq{}}\PY{l+s+s2}{amsterdam}\PY{l+s+s2}{\PYZdq{}}\PY{p}{,} \PY{l+s+s2}{\PYZdq{}}\PY{l+s+s2}{tokyo}\PY{l+s+s2}{\PYZdq{}}\PY{p}{,} \PY{l+s+s2}{\PYZdq{}}\PY{l+s+s2}{paris}\PY{l+s+s2}{\PYZdq{}}\PY{p}{]}\PY{p}{)}\PY{p}{)}
        \end{Verbatim}
    \end{tcolorbox}

    \begin{Verbatim}[commandchars=\\\{\}]
[1 0 3 2]
    \end{Verbatim}

    \hypertarget{spliting-the-dataset-into-training-set-and-test-set}{%
        \subsection{Spliting the dataset into Training set and Test
        set}\label{spliting-the-dataset-into-training-set-and-test-set}}

    \begin{tcolorbox}[breakable, size=fbox, boxrule=1pt, pad at break*=1mm,colback=cellbackground, colframe=cellborder]
        \prompt{In}{incolor}{40}{\boxspacing}
        \begin{Verbatim}[commandchars=\\\{\}]
\PY{k+kn}{from}\PY{+w}{ }\PY{n+nn}{sklearn}\PY{n+nn}{.}\PY{n+nn}{model\PYZus{}selection}\PY{+w}{ }\PY{k+kn}{import} \PY{n}{train\PYZus{}test\PYZus{}split}
\PY{n}{X\PYZus{}train}\PY{p}{,} \PY{n}{X\PYZus{}test}\PY{p}{,} \PY{n}{y\PYZus{}train}\PY{p}{,} \PY{n}{y\PYZus{}test} \PY{o}{=} \PY{n}{train\PYZus{}test\PYZus{}split}\PY{p}{(}\PY{n}{X}\PY{p}{,} \PY{n}{y}\PY{p}{,} \PY{n}{test\PYZus{}size}\PY{o}{=}\PY{l+m+mf}{0.2}\PY{p}{,} \PY{n}{random\PYZus{}state}\PY{o}{=}\PY{l+m+mi}{1}\PY{p}{)}    \PY{c+c1}{\PYZsh{} random\PYZhy{}state is used for reproducibility}
        \end{Verbatim}
    \end{tcolorbox}

    \begin{tcolorbox}[breakable, size=fbox, boxrule=1pt, pad at break*=1mm,colback=cellbackground, colframe=cellborder]
        \prompt{In}{incolor}{41}{\boxspacing}
        \begin{Verbatim}[commandchars=\\\{\}]
\PY{n+nb}{print}\PY{p}{(}\PY{n}{X\PYZus{}train}\PY{p}{)}
        \end{Verbatim}
    \end{tcolorbox}

    \begin{Verbatim}[commandchars=\\\{\}]
[[0.0 0.0 1.0 38.77777777777778 52000.0]
 [0.0 1.0 0.0 40.0 63777.77777777778]
 [1.0 0.0 0.0 44.0 72000.0]
 [0.0 0.0 1.0 38.0 61000.0]
 [0.0 0.0 1.0 27.0 48000.0]
 [1.0 0.0 0.0 48.0 79000.0]
 [0.0 1.0 0.0 50.0 83000.0]
 [1.0 0.0 0.0 35.0 58000.0]]
    \end{Verbatim}

    \begin{tcolorbox}[breakable, size=fbox, boxrule=1pt, pad at break*=1mm,colback=cellbackground, colframe=cellborder]
        \prompt{In}{incolor}{42}{\boxspacing}
        \begin{Verbatim}[commandchars=\\\{\}]
\PY{n+nb}{print}\PY{p}{(}\PY{n}{X\PYZus{}test}\PY{p}{)}
        \end{Verbatim}
    \end{tcolorbox}

    \begin{Verbatim}[commandchars=\\\{\}]
[[0.0 1.0 0.0 30.0 54000.0]
 [1.0 0.0 0.0 37.0 67000.0]]
    \end{Verbatim}

    \begin{tcolorbox}[breakable, size=fbox, boxrule=1pt, pad at break*=1mm,colback=cellbackground, colframe=cellborder]
        \prompt{In}{incolor}{43}{\boxspacing}
        \begin{Verbatim}[commandchars=\\\{\}]
\PY{n+nb}{print}\PY{p}{(}\PY{n}{y\PYZus{}train}\PY{p}{)}
        \end{Verbatim}
    \end{tcolorbox}

    \begin{Verbatim}[commandchars=\\\{\}]
[0 1 0 0 1 1 0 1]
    \end{Verbatim}

    \begin{tcolorbox}[breakable, size=fbox, boxrule=1pt, pad at break*=1mm,colback=cellbackground, colframe=cellborder]
        \prompt{In}{incolor}{44}{\boxspacing}
        \begin{Verbatim}[commandchars=\\\{\}]
\PY{n+nb}{print}\PY{p}{(}\PY{n}{y\PYZus{}test}\PY{p}{)}
        \end{Verbatim}
    \end{tcolorbox}

    \begin{Verbatim}[commandchars=\\\{\}]
[0 1]
    \end{Verbatim}

    \hypertarget{features-scaling}{%
        \subsection{Features Scaling}\label{features-scaling}}

    \begin{tcolorbox}[breakable, size=fbox, boxrule=1pt, pad at break*=1mm,colback=cellbackground, colframe=cellborder]
        \prompt{In}{incolor}{45}{\boxspacing}
        \begin{Verbatim}[commandchars=\\\{\}]
\PY{k+kn}{from}\PY{+w}{ }\PY{n+nn}{sklearn}\PY{n+nn}{.}\PY{n+nn}{preprocessing}\PY{+w}{ }\PY{k+kn}{import} \PY{n}{StandardScaler}    \PY{c+c1}{\PYZsh{} Z\PYZus{}score = [(x \PYZhy{} x\PYZus{}bar) / s] or [(x \PYZhy{} mu) / sigma]}
\PY{n}{sc} \PY{o}{=} \PY{n}{StandardScaler}\PY{p}{(}\PY{p}{)}                               \PY{c+c1}{\PYZsh{} x \PYZhy{}\PYZgt{} data}
\PY{n}{X\PYZus{}train}\PY{p}{[}\PY{p}{:}\PY{p}{,} \PY{l+m+mi}{3}\PY{p}{:}\PY{p}{]} \PY{o}{=} \PY{n}{sc}\PY{o}{.}\PY{n}{fit\PYZus{}transform}\PY{p}{(}\PY{n}{X\PYZus{}train}\PY{p}{[}\PY{p}{:}\PY{p}{,} \PY{l+m+mi}{3}\PY{p}{:}\PY{p}{]}\PY{p}{)}   \PY{c+c1}{\PYZsh{} x\PYZus{}bar / mu \PYZhy{}\PYZgt{} sample / population mean}
\PY{n}{X\PYZus{}test}\PY{p}{[}\PY{p}{:}\PY{p}{,} \PY{l+m+mi}{3}\PY{p}{:}\PY{p}{]} \PY{o}{=} \PY{n}{sc}\PY{o}{.}\PY{n}{transform}\PY{p}{(}\PY{n}{X\PYZus{}test}\PY{p}{[}\PY{p}{:}\PY{p}{,} \PY{l+m+mi}{3}\PY{p}{:}\PY{p}{]}\PY{p}{)}         \PY{c+c1}{\PYZsh{} s / sigma \PYZhy{}\PYZgt{} sample / population  standard deviation}
        \end{Verbatim}
    \end{tcolorbox}

    \begin{tcolorbox}[breakable, size=fbox, boxrule=1pt, pad at break*=1mm,colback=cellbackground, colframe=cellborder]
        \prompt{In}{incolor}{46}{\boxspacing}
        \begin{Verbatim}[commandchars=\\\{\}]
\PY{n+nb}{print}\PY{p}{(}\PY{n}{X\PYZus{}train}\PY{p}{)}
        \end{Verbatim}
    \end{tcolorbox}

    \begin{Verbatim}[commandchars=\\\{\}]
[[0.0 0.0 1.0 -0.19159184384578545 -1.0781259408412425]
 [0.0 1.0 0.0 -0.014117293757057777 -0.07013167641635372]
 [1.0 0.0 0.0 0.566708506533324 0.633562432710455]
 [0.0 0.0 1.0 -0.30453019390224867 -0.30786617274297867]
 [0.0 0.0 1.0 -1.9018011447007988 -1.420463615551582]
 [1.0 0.0 0.0 1.1475343068237058 1.232653363453549]
 [0.0 1.0 0.0 1.4379472069688968 1.5749910381638885]
 [1.0 0.0 0.0 -0.7401495441200351 -0.5646194287757332]]
    \end{Verbatim}

    \begin{tcolorbox}[breakable, size=fbox, boxrule=1pt, pad at break*=1mm,colback=cellbackground, colframe=cellborder]
        \prompt{In}{incolor}{47}{\boxspacing}
        \begin{Verbatim}[commandchars=\\\{\}]
\PY{n+nb}{print}\PY{p}{(}\PY{n}{X\PYZus{}test}\PY{p}{)}
        \end{Verbatim}
    \end{tcolorbox}

    \begin{Verbatim}[commandchars=\\\{\}]
[[0.0 1.0 0.0 -1.4661817944830124 -0.9069571034860727]
 [1.0 0.0 0.0 -0.44973664397484414 0.2056403393225306]]
    \end{Verbatim}


    % Add a bibliography block to the postdoc



\end{document}
